\chapter{Experimento Preliminar}

Este capítulo presenta la estrategia propuesta para validar el funcionamiento del prototipo integrado. 
Para ello, se evaluarán los resultados de las pruebas funcionales individuales realizadas durante 
el desarrollo de los subsistemas y se ejecutarán pruebas de integración del sistema completo en un 
laboratorio de electrónica. La evaluación integral permitirá comprobar la ejecución del ciclo de 
prueba y determinar la capacidad del prototipo para generar registros 
trazables de las verificaciones de pulseras antiestáticas.

La validación está planificada en el laboratorio TEM de Intel Costa Rica, debido
a que el proyecto surge de una necesidad identificada en ese entorno y las
condiciones de operación corresponden con su aplicación prevista. Si
circunstancias fuera del control del proyecto impiden realizar las pruebas en
TEM, se utilizará como alternativa un laboratorio de la Escuela de Ingeniería
Electrónica del Tecnológico de Costa Rica, Campus Tecnológico Local San Carlos.
En ambos casos se aplicarán el mismo procedimiento, los mismos indicadores y los
mismos criterios de aceptación.

\section{Condiciones del experimento}

Antes de ejecutar el experimento se realizarán pruebas funcionales de los
subsistemas principales. Se comprobará la lectura de los estados
\textit{LO}, \textit{OK} y \textit{HI} entregados por el probador HZR-171, la
identificación mediante código de barras, el reconocimiento facial
complementario, la navegación de la interfaz gráfica y la escritura de datos en
la memoria local. 

Para la evaluación integral se ejecutarán al menos 30 ciclos completos con un
mínimo de cinco usuarios de prueba. Los identificadores utilizados podrán ser
ficticios o anonimizados. El reconocimiento facial será opcional y únicamente
se probará con participantes que autoricen previamente su uso.

Un ciclo completo estará formado por las siguientes etapas:

\begin{enumerate}
\item Identificación del usuario mediante código de barras o, de forma
complementaria, reconocimiento facial.
\item Presentación de las instrucciones en la interfaz gráfica.
\item Ejecución de la prueba de pulsera antiestática.
\item Adquisición e interpretación del resultado entregado por el HZR-171.
\item Almacenamiento local del registro y presentación del resultado al usuario.
\end{enumerate}

Cada registro deberá contener, como mínimo, el nombre y el identificador del
usuario, la fecha, la hora y el resultado de la prueba ESD. Además, incorporará la
temperatura y la humedad relativa disponibles durante el ciclo. 

\section{Indicadores de validación}

La evaluación del prototipo se concentrará en dos indicadores: el índice de
trazabilidad \(TZ\) y la tasa de integración funcional \(IF\). De esta manera se
evita utilizar varios porcentajes que evalúan partes de un mismo ciclo y se
mantiene el análisis relacionado con los objetivos principales del proyecto.

\subsection{Índice de trazabilidad}

El índice de trazabilidad representa el porcentaje de ciclos ejecutados que
generan un registro completo:

\begin{equation}
TZ =
\frac{N_{\text{registros completos}}}
{N_{\text{ciclos ejecutados}}}
\times 100\%
\label{eq:indice_trazabilidad}
\end{equation}

Para este indicador, un registro completo es aquel que permite relacionar de
forma inequívoca al usuario con la fecha, la hora y el resultado ESD obtenido.
La temperatura y la humedad relativa se almacenarán como datos complementarios,
pero no determinarán por sí solas si el registro es trazable. El criterio de
aceptación será \(TZ \geq 95\%\).

\subsection{Tasa de integración funcional}

La tasa de integración funcional representa el porcentaje de ciclos en los que
la secuencia completa se ejecuta correctamente:

\begin{equation}
IF =
\frac{N_{\text{ciclos funcionalmente correctos}}}
{N_{\text{ciclos ejecutados}}}
\times 100\%
\label{eq:integracion_funcional}
\end{equation}

Un ciclo será funcionalmente correcto cuando el sistema identifique al usuario,
guíe la operación mediante la interfaz, interprete el estado del probador,
muestre el resultado correspondiente y almacene el registro sin errores que
impidan completar el flujo. El criterio de aceptación será \(IF \geq 95\%\).

\section{Procedimiento experimental}

El experimento se realizará siguiendo este procedimiento:

\begin{enumerate}

\item Ejecutar al menos 30 ciclos completos, alternando los usuarios y los
resultados de prueba disponibles.

\item Revisar cada registro almacenado y comprobar que el nombre y el
identificador del usuario, la fecha, la hora, el resultado ESD, la temperatura
y la humedad relativa coincidan con los datos observados durante el ciclo.

\item Utilizar la interfaz gráfica del prototipo para consultar los registros y
generar un reporte histórico en formato CSV correspondiente al conjunto de
ciclos ejecutados.

\item Verificar que el reporte generado coincida con los registros almacenados,
incluya los campos definidos y presente la cantidad de registros completos
obtenidos durante los 30 ciclos.

\item Contabilizar los registros completos y los ciclos funcionalmente
correctos.

\item Calcular los indicadores \(TZ\) e \(IF\) y compararlos con sus criterios de
aceptación.

\end{enumerate}

Durante las pruebas también se documentarán errores de comunicación, fallas de
lectura, problemas de navegación y cualquier ajuste necesario en el montaje
físico. Las observaciones cualitativas servirán para complementar los indicadores
y orientar correcciones del diseño.

\section{Criterios de aceptación}

La Tabla~\ref{tab:criterios_validacion} resume los criterios principales del
experimento.

\begin{table}[H]
\centering
\caption{Criterios de aceptación del experimento preliminar.}
\label{tab:criterios_validacion}
\begin{tabularx}{\textwidth}{>{\raggedright\arraybackslash}p{3.4cm}
                            >{\raggedright\arraybackslash}X
                            >{\raggedright\arraybackslash}p{3.5cm}}
\toprule
\textbf{Evaluación} &
\textbf{Evidencia} &
\textbf{Criterio de aceptación} \\
\midrule
Trazabilidad \(TZ\) &
Registros que asocian usuario, fecha, hora y resultado ESD respecto al total de
ciclos ejecutados. &
\(TZ \geq 95\%\) en al menos 30 ciclos. \\
\addlinespace
Integración funcional \(IF\) &
Ciclos completados correctamente desde la identificación hasta el
almacenamiento del registro. &
\(IF \geq 95\%\) en al menos 30 ciclos. \\
\addlinespace
Reporte histórico &
Archivo CSV generado a partir de los registros almacenados. &
Al menos un reporte legible y coherente con los ciclos ejecutados. \\
\bottomrule
\end{tabularx}
\end{table}

\section{Análisis de resultados}

El prototipo se considerará satisfactorio si cumple los criterios establecidos para \(TZ\) e \(IF\).
El análisis presentará los valores obtenidos, las fallas observadas y las
correcciones realizadas durante la integración. También se informará el
laboratorio donde finalmente se ejecutó el experimento y, si fue necesario
utilizar la alternativa de TEC San Carlos, se explicará la circunstancia externa
que impidió realizar las pruebas en TEM.

Los resultados se presentarán de forma agregada o anonimizada. No se publicarán
nombres, fotografías ni identificaciones empresariales de los participantes. El
experimento demostrará el desempeño funcional del prototipo dentro del alcance
académico definido, pero no constituirá una certificación industrial del
sistema.
